\section{Введение}

Денежно-кредитная политика проводится не по истинному состоянию экономики, а по неполному и
шумному представлению о нём. Центральный банк располагает макроэкономическими индикаторами
инфляции, выпуска, потребительской активности и финансовых условий, однако эти показатели
наблюдаются с ошибками измерения, задержками и последующими пересмотрами. При этом распределение
ликвидности, активов и обязательств домохозяйств не наблюдается напрямую. Литература о данных в
реальном времени (\textit{real-time data}), начиная с Orphanides
\cite{orphanides2001realtime}, показывает, что рекомендации правил денежно-кредитной политики
могут существенно зависеть от доступной на момент решения информации.

В HANK-моделях это ограничение особенно существенно. Реакция потребления на изменение ставки
зависит от распределения домохозяйств по активам, доходам, ликвидности и предельной склонности к
потреблению. Kaplan, Moll и Violante \cite{kaplan2018hank} показывают, что неоднородность
домохозяйств меняет силу денежно-кредитной трансмиссии через различия в ликвидности и предельной
склонности к потреблению (\textit{marginal propensity to consume}, MPC). Auclert
\cite{auclert2019redistribution} выделяет перераспределительные каналы денежно-кредитной
политики, включая различия в предельной склонности к потреблению, канал Фишера
(\textit{Fisher channel}) и чувствительность процентных доходов и расходов к ставке
(\textit{interest rate exposure}). Поэтому агрегатные показатели могут быть недостаточны для
описания будущей реакции потребления, выпуска и инфляции на изменение ставки.

Однако из HANK-мотивации не следует автоматическая полезность любых распределительных данных.
Распределительные статистики имеют содержательную ценность только в той мере, в какой они
улучшают оценку состояния экономики, важного для будущей реакции потребления, выпуска и инфляции
на изменение ставки. В работе рассматривается вопрос о предельной ценности распределительной
информации сверх агрегатной оценки состояния, уже построенной по шумным макроэкономическим
наблюдениям.

Цель работы состоит в том, чтобы оценить предельную ценность распределительных сигналов для
правила процентной ставки при неполной наблюдаемости. Для достижения цели решаются четыре задачи.
Сначала HANK/SSJ-отклики переводятся в низкоразмерную модель в пространстве состояний. Затем
задаются несколько информационных режимов регулятора, различающихся доступными наблюдениями и
способом их сжатия. Далее для каждого режима заново настраивается однотипное правило процентной
ставки. Наконец, полученные правила сравниваются по значениям одной и той же функции потерь
стабилизации на общей тестовой выборке.

Работа устроена следующим образом. В разделе 2 описывается экономическая среда и распределительные
статистики. В разделе 3 вводятся неполная наблюдаемость и информационные режимы регулятора. В
разделе 4 формулируется задача сравнения информации, правило процентной ставки и критерий
потерь. Раздел 5 обсуждает экономический механизм распределительных сигналов. В разделе 6
описан дизайн вычислительного эксперимента. В разделе 7 представлены результаты, в разделе 8
обсуждается их интерпретация, а раздел 9 содержит заключение.
